%! Parent Functions and Transformations — Unit 2, Lesson 2.2 — Group Activity
\documentclass[10pt]{article}
\usepackage{atda-article}
\usepackage{atda-boxes}
\newcommand{\TallMath}[1]{$\displaystyle #1\rule[-1.4em]{0pt}{3.2em}$}
\setlength{\workrowsep}{6pt}

\begin{document}

\pageheader{Unit 2, Lesson 2.2}{Group Activity --- Moved, Flipped, Stretched}
% NAMESTRIP: no \namedateperiod here. The student writes their name once, on the
% cover this component is stapled behind. See references/conventions.md.

\vspace{0.15in}

\begin{scenariobox}[The Scenario]{royal}
\small
Every graph below is one of the three parents with something done to it. Your group's job is the
same every time: \textbf{name the parent, name what was done to it, and say which of the domain and
the range moved.} Work together. \textbf{Everyone starts at Tier R.} When your whole group can do
Tier R without help, move on.
\end{scenariobox}

\vspace{0.10in}

\boxguard[30]
\begin{tcolorbox}[colback=white, colframe=black!40, title=\textbf{Tier R --- Remediate},
    fonttitle=\bfseries, arc=2mm, left=3mm, right=3mm, top=2mm, bottom=2mm, breakable]
\small
The quadratic parent $f(x)=x^2$ is drawn with two of its children.

\begin{center}
\begin{tikzpicture}
\begin{axis}[
    width=10.4cm, height=4.8cm,
    xlabel={$x$}, ylabel={$y$},
    xmin=-6.6, xmax=3.6, ymin=-6, ymax=10.5,
    xtick={-6,-5,-4,-3,-2,-1,0,1,2,3}, ytick={-5,0,5,10},
    grid=both, grid style={linegray!45}, axis lines=left,
    tick label style={font=\scriptsize}, label style={font=\scriptsize},
    legend entries={{$y=x^2$},{$y=x^2-5$},{$y=(x+3)^2$}},
    legend style={font=\scriptsize, draw=none, fill=none,
      at={(0.5,1.02)}, anchor=south, legend columns=3}]
  \addplot[royal, very thick, domain=-3:3, samples=80]{x^2};
  \addplot[goldacc, very thick, dashed, domain=-3:3, samples=80]{x^2-5};
  \addplot[plumacc, very thick, dash dot, domain=-6:0, samples=80]{(x+3)^2};
\end{axis}
\end{tikzpicture}
\end{center}

\begin{enumerate}[leftmargin=1.6em, itemsep=4pt, topsep=2pt]
  \item For each child, say which way and how far it moved, then write it in function notation
        starting from $f(x)=x^2$.

\writelines{2}

  \item State the range of all three. Which child changed the range, and how does the position of
        its number (inside or outside the parentheses) explain that?

\writelines{2}

  \item No graph this time. Name the parent, then describe the move.
        \quad (a) $y=x+7$ \quad (b) $y=(x-6)^2$ \quad (c) $y=2^x-1$

\writelines{3}

  \item In one sentence: why does $y=(x+3)^2$ move the graph \emph{left}, even though the equation
        shows a plus sign?

\writelines{1}
\end{enumerate}
\end{tcolorbox}

\vspace{0.10in}

\boxguard[30]
\begin{tcolorbox}[colback=white, colframe=black!40,
    title=\textbf{Tier A --- Approaching Proficiency},
    fonttitle=\bfseries, arc=2mm, left=3mm, right=3mm, top=2mm, bottom=2mm, breakable]
\small
Same parent, but nothing here is a slide.

\begin{center}
\begin{tikzpicture}
\begin{axis}[
    width=10.0cm, height=4.8cm,
    xlabel={$x$}, ylabel={$y$},
    xmin=-3.6, xmax=3.6, ymin=-10, ymax=10,
    xtick={-3,-2,-1,0,1,2,3}, ytick={-9,-6,-3,0,3,6,9},
    grid=both, grid style={linegray!45}, axis lines=left,
    tick label style={font=\scriptsize}, label style={font=\scriptsize},
    legend entries={{$y=x^2$},{$y=-x^2$},{$y=\tfrac{1}{3}x^2$}},
    legend style={font=\scriptsize, draw=none, fill=none,
      at={(0.5,1.02)}, anchor=south, legend columns=3}]
  \addplot[royal, very thick, domain=-3:3, samples=80]{x^2};
  \addplot[goldacc, very thick, dashed, domain=-3:3, samples=80]{-x^2};
  \addplot[plumacc, very thick, dash dot, domain=-3:3, samples=80]{x^2/3};
\end{axis}
\end{tikzpicture}
\end{center}

\begin{enumerate}[leftmargin=1.6em, itemsep=4pt, topsep=2pt]
  \item Describe each child as a transformation of $f(x)=x^2$, in words and in function notation.

\writelines{2}

  \item State the range of all three. Two of them share a range. Explain why the dilation left the
        range alone while the reflection did not.

\writelines{2}

  \item \textbf{In context.} A conveyor belt repeats one cycle, and $C(t)$ is the height of the
        lifting arm $t$ seconds into the cycle, for $0 \le t \le 8$. The factory doubles the belt
        speed, so the arm now does the identical motion in half the time. The new function is
        $C(2t)$.
        \begin{enumerate}[label=(\alph*), leftmargin=1.6em, itemsep=3pt, topsep=3pt]
          \item What is the domain of the new function, and what is its range?

\writelines{1}

          \item Is this a horizontal stretch or a horizontal compression? Point to the value of
                $|k|$ that decides it.

\writelines{1}

          \item Which moved --- the domain or the range --- and why is that the answer you should
                have predicted before doing anything?

\writelines{1}
        \end{enumerate}

  \item The exponential parent is $h(x)=2^x$, which climbs. Its reflection $h(-x)=2^{-x}$ falls.
        Use the table below to say why, then name the axis it was flipped over.
        \smallskip

        \renewcommand{\arraystretch}{1.3}
        \begin{tabular}{@{}|c|c|c|c|c|c|@{}}
        \hline
        $x$ & $-2$ & $-1$ & $0$ & $1$ & $2$ \\ \hline
        $2^{-x}$\rule[-0.75em]{0pt}{1.9em} & \blank{1.1cm} & \blank{1.1cm} & \blank{1.1cm} &
            \blank{1.1cm} & \blank{1.1cm} \\ \hline
        \end{tabular}
        \smallskip

\writelines{1}
\end{enumerate}
\end{tcolorbox}

\vspace{0.10in}

\boxguard[30]
\begin{tcolorbox}[colback=white, colframe=black!40, title=\textbf{Tier E --- Extension},
    fonttitle=\bfseries, arc=2mm, left=3mm, right=3mm, top=2mm, bottom=2mm, breakable]
\small

\begin{enumerate}[leftmargin=1.6em, itemsep=4pt, topsep=2pt]
  \item \textbf{Critique the work.} Each student below is wrong. Name the error and give the
        correct answer.
        \begin{enumerate}[label=(\alph*), leftmargin=1.6em, itemsep=3pt, topsep=3pt]
          \item ``$y=(x+5)^2$ is the parabola moved \emph{right} $5$, because the $5$ is
                positive.''

\writelines{2}

          \item ``$y=-x^2$ and $y=(-x)^2$ are the same graph. A minus sign is a minus sign.''

\writelines{2}
        \end{enumerate}

  \item \textbf{Two moves at once.} Consider $y=(x-3)^2+2$.
        \begin{enumerate}[label=(\alph*), leftmargin=1.6em, itemsep=3pt, topsep=3pt]
          \item Name the parent, then name \emph{both} transformations, saying which number is
                inside the parentheses and which is outside.

\writelines{2}

          \item Where is the turning point of the new graph? Check your answer by evaluating.
\begin{work}
  y &= (3-3)^2+2 \\
    &= 0+2 \\
    &= 2
\end{work}

          \item State the domain and the range of $y=(x-3)^2+2$ in interval notation. Which one
                did the $+2$ move, and which one did the $-3$ move?

\writelines{2}
        \end{enumerate}

  \item \textbf{Justify.} A classmate says: ``The graph of $g(x)=x^2$ can't ever have been
        reflected over the $y$-axis --- it looks exactly the same afterwards, so nothing
        happened.'' Are they right? Use the meaning of $f(-x)$ to answer, then give a parent
        function for which the same flip obviously \emph{does} change the picture.

\writelines{3}
\end{enumerate}
\end{tcolorbox}

\end{document}
